% Generated by Listas/make_statements.py from pratica-lecture-05.md, item 16. Do not edit.
\textbf{Source.} Gibbons (1992), \emph{Game Theory for Applied Economists}, §2.3 Problems, Problem 2.11, p.~135 {[}PDF 142{]}. Answer in the source: none in the book (Muñoz-Garcia Ch. 6 Ex. 7, p.~299, works the same game). Reproduced verbatim. (The same game is Muñoz-Garcia \& Toro-Gonzalez Ch. 6 Ex. 7, p.~299, with a worked answer.)



\begin{quote}

2.11. The simultaneous-move game (below) is played twice, with the outcome of the first stage observed before the second stage begins. There is no discounting. Can the payoff \((4,4)\) be achieved in the first stage in a pure-strategy subgame-perfect Nash equilibrium? If so, give strategies that do so. If not, prove why not.



{\def\LTcaptype{none} % do not increment counter

\begin{longtable}[]{@{}llll@{}}

\toprule\noalign{}

& \(L\) & \(C\) & \(R\) \\

\midrule\noalign{}

\endhead

\bottomrule\noalign{}

\endlastfoot

\(T\) & 3, 1 & 0, 0 & 5, 0 \\

\(M\) & 2, 1 & 1, 2 & 3, 1 \\

\(B\) & 1, 2 & 0, 1 & 4, 4 \\

\end{longtable}

}

\end{quote}

